\begin{abstract}
This thesis reports work carried out during a 2026 internship with the MILeS group at LAMSADE. It asks how much task-relevant information must cross a constrained representation to preserve learned behavior.

The first part treats a low-rank fine-tuning adapter as a message whose receiver already has the frozen base model. Rate is the exact size of the serialized adapter after reload, rather than nominal rank or bit width. Distortion is measured by held-out prediction gain and exact task behavior. The strongest current result is a screened model-relative measure: the spectral log-volume of the output-head corrections requested by the dataset has mean within-model Spearman correlation 0.764 with the smallest useful adapter rate across 22 distinct model--dataset arms. It reduces leave-one-task-family-out prediction error by 29\% against a model-only fit, while base-model code length gives no useful reduction. The ordering of five corpora changes between Mistral-7B and Qwen2.5-7B, and fixed-row diversity raises the threshold on XBRL for every seed and both models but not reliably on SQL. These results support information in the learned dataset relative to the frozen receiver as the missing quantity. They screen a fixed gradient sketch on two models rather than establish a scaling law.

The second part studies reasoning state. Token partitions for answer alignment, symbolic updates, correctness change, and trajectory compression differ across more than 370,000 held-out tokens, even though supervised probes recover each ontology. An exact finite controller gives two representations identical initial mutual information but sharply different long-horizon regret because one forgets future observations. Language-model interventions show that a canonical four-bit proof-state channel can remain exact through horizon 256 and reach 97.22\% final accuracy, compared with 72.22\% for matched one-pass training. Lossy and path-aliased codes fail in different ways. Local register updates that are roughly 89\% reliable still compound to near-chance answers at horizon 32. These results support an exact, closed-loop, utility-relative account of information in learning and reasoning. They do not establish a universal information law or a canonical unit of thought.
\end{abstract}
